\documentclass[12pt, a4paper]{article}

% --- Required Packages ---
\usepackage[utf8]{inputenc}
\usepackage{amsmath, amssymb, amsthm}
\usepackage{geometry}
\geometry{top=2.5cm, bottom=2.5cm, left=2.5cm, right=2.5cm}
\usepackage{enumitem}

\newtheorem{theorem}{Theorem}
\newtheorem{lemma}{Lemma}
\newtheorem{proposition}{Proposition}

\begin{document}

\subsection*{B. Algorithm and Bicriteria Bound for the MAL-C Problem without a Time-Expanded Graph}

The MAL-C problem (Definition~3 of the project plan) requires us to choose, for every edge $e\in E$, a label set $\lambda(e)\subseteq\{1,\dots,a\}$, and for every demand $d_i=(\mathrm{src}(d_i),\mathrm{dst}(d_i),\mathrm{rel}(d_i),\mathrm{dem}(d_i))$ a flow $f_i(e,t)$, so as to minimise
\[
\mathrm{cost}(\lambda,\{f_i\}) \;=\; \sum_{e\in E}\sum_{t=1}^{a}\Bigl\lceil \textstyle\sum_{i=1}^{h} f_i(e,t)\Bigr\rceil .
\]
The standard route to a solution is to construct a Time-Expanded Graph (TEG) of size $\Theta(a\cdot|V|+a\cdot|E|)$. When $a$ is much larger than $\mathrm{Diam}(G)$ this becomes pseudo-polynomial.  We avoid the TEG entirely and obtain a \emph{bicriteria} bound (multiplicative + additive) whose multiplicative factor depends only on $\mathrm{Diam}(G)$.

\paragraph{Notation.} $n=|V|$, $m=|E|$. Edge weights $w(e)\in\mathbb N$ are travel times. All graph distances $\mathrm{dist}_G(\cdot,\cdot)$ are measured in \emph{hops} (number of edges) unless explicitly written $w$-distance.

\subsubsection*{1. Algorithm: Hub-Routed Per-Edge Schedule (HRPES)}

We separate \emph{spatial routing} (which edges are traversed) from \emph{temporal scheduling} (at which times those edges carry flow), and we attach a per-edge label rather than a per-node time.

\paragraph{Phase 1 (spatial backbone).}
Apply Theorem~1 of \cite{CDO25} with $\delta=\tfrac12$ and $n_2=\sqrt{n\log n}$ to obtain a $\bigl\lfloor 0.5\,\mathrm{Diam}(G)\bigr\rfloor$-Dominating Set $S\subseteq V$ with $|S|=O(\sqrt{n\log n})$.
\begin{itemize}[noitemsep]
  \item For every $v\in V$ fix a unique \emph{home hub} $\mathrm{hub}(v)\in S$ minimising hop-distance from $v$ to $S$. Compute a hop-shortest path $\pi_v$ from $v$ to $\mathrm{hub}(v)$. By choice of $d=\lfloor 0.5\,\mathrm{Diam}(G)\rfloor$, $|\pi_v|\le d$.
  \item For every ordered pair $(s,s')\in S\times S$ compute a hop-shortest path $\pi_{s,s'}$ in $G$, of length $\le \mathrm{Diam}(G)$.
\end{itemize}
The backbone is $H_\mathrm{hub} = \bigl(\bigcup_v \pi_v\bigr)\cup\bigl(\bigcup_{s,s'\in S}\pi_{s,s'}\bigr)$.
For each demand $d_i$ the algorithm uses the path
\[
P_i:\quad \mathrm{src}(d_i)\;\xrightarrow{\pi_{\mathrm{src}(d_i)}}\;\mathrm{hub}(\mathrm{src}(d_i))\;\xrightarrow{\pi_{\mathrm{hub}(\mathrm{src}(d_i)),\mathrm{hub}(\mathrm{dst}(d_i))}}\;\mathrm{hub}(\mathrm{dst}(d_i))\;\xrightarrow{\pi_{\mathrm{dst}(d_i)}^{-1}}\;\mathrm{dst}(d_i),
\]
whose hop-length satisfies $k_i^{\mathrm{algo}} = |P_i| \le d+\mathrm{Diam}(G)+d \le 2\,\mathrm{Diam}(G)$.

\paragraph{Phase 2 (per-edge schedule, computed for each demand).}
For each demand $d_i$, walk $P_i$ \emph{forward} from $\mathrm{src}(d_i)$ and assign to its $j$-th edge $e_{i,j}=(u_{j-1},u_j)$ a departure time
\[
\tau_i(e_{i,j}) \;=\; \mathrm{rel}(d_i)+\sum_{\ell<j} w(e_{i,\ell}).
\]
This corresponds to the demand departing at $\mathrm{rel}(d_i)$ and using each edge as soon as it arrives at the previous node (no waiting).

\emph{Feasibility check.} We must have $\tau_i(e_{i,k_i^{\mathrm{algo}}})+w(e_{i,k_i^{\mathrm{algo}}})\le a$, i.e.\ the total $w$-length of $P_i$ is at most $a-\mathrm{rel}(d_i)$.
We \emph{define} a demand to be \emph{schedulable on the backbone} if this condition holds; otherwise it is \emph{escapable}, meaning even the original-graph optimal $w$-shortest path may be too long, in which case the demand is infeasible and a feasible MAL-C instance cannot exist for it (this is an instance-level check, not a phase-3 fallback).

\paragraph{Phase 3 (label and flow assignment).}
Set
\[
\lambda(e) \;=\; \bigl\{\tau_i(e) \;:\; e\in P_i\bigr\}\subseteq\{1,\dots,a\},\qquad
f_i(e,t) \;=\; \begin{cases}\mathrm{dem}(d_i) & \text{if } e\in P_i \text{ and } t=\tau_i(e),\\ 0 & \text{otherwise.}\end{cases}
\]
Flow conservation, label--flow consistency, and temporal connectivity are immediate from the construction.

\paragraph{Optional consolidation step.}
If two demands $d_i,d_j$ share an edge $e$ in their backbone paths and $\tau_i(e)\ne\tau_j(e)$, but both have slack $a-\mathrm{rel}-\sum w$ on $P_i,P_j$, we may shift one of them to use the same time slot, paying the difference in waiting time at intermediate nodes (free, since node storage is unbounded). This step only ever \emph{decreases} the cost; the worst-case bound below is proven without using it.

\subsubsection*{2. Lower Bound (Fluid Relaxation)}

\begin{lemma}[Fluid lower bound]
\label{lem:fluid}
Let $\mathrm{DP}[v][k]$ be the earliest arrival time at $v$ after exactly $k$ edge traversals starting from $\mathrm{src}(d_i)$ at time $\mathrm{rel}(d_i)$:
$\mathrm{DP}[\mathrm{src}(d_i)][0]=\mathrm{rel}(d_i)$, and
\[
\mathrm{DP}[v][k] = \min_{(u,v)\in E}\{\mathrm{DP}[u][k-1]+w(u,v)\}.
\]
Define $k_i^* = \min\{k\;:\;\mathrm{DP}[\mathrm{dst}(d_i)][k]\le a\}$. Then
\[
\mathrm{OPT}_\mathrm{fluid} \;:=\; \sum_{i=1}^{h}\mathrm{dem}(d_i)\cdot k_i^* \;\le\; \mathrm{OPT}_\mathrm{MAL\text{-}C}.
\]
\end{lemma}
\begin{proof}
Drop the ceiling in the MAL-C objective and call the relaxed minimum $\mathrm{OPT}_\mathrm{fluid}$. Since $\lceil x\rceil\ge x$ for $x\ge 0$, we have $\mathrm{OPT}_\mathrm{fluid}\le \mathrm{OPT}_\mathrm{MAL\text{-}C}$. Because node storage is free and labels are unconstrained, in the relaxed problem each demand independently picks a path minimising its hop count subject to the time-feasibility $\mathrm{DP}[\mathrm{dst}(d_i)][k]\le a$, giving exactly $\mathrm{dem}(d_i)\cdot k_i^*$ contribution per demand and $\mathrm{OPT}_\mathrm{fluid}=\sum_i \mathrm{dem}(d_i)\cdot k_i^*$.
\end{proof}

\subsubsection*{3. Cost of HRPES (Bicriteria Bound)}

\paragraph{Hop stretch.}
By construction $k_i^{\mathrm{algo}}\le k_i^*+2d \le k_i^*+\mathrm{Diam}(G)$, and also $k_i^{\mathrm{algo}}\le 2\,\mathrm{Diam}(G)$. Using $k_i^*\ge 1$ gives the multiplicative form $k_i^{\mathrm{algo}}\le 2\,\mathrm{Diam}(G)\cdot k_i^*$, but this is loose; the additive form is tight up to a constant.

\paragraph{Counting used edge--time pairs.}
Let $U=\bigl\{(e,t):e\in P_i,\;t=\tau_i(e),\;\text{some }i\bigr\}$ be the set of edge--time pairs that carry positive flow. Each pair contributes a ceiling overhead of at most $1-\bigl(\sum_i f_i(e,t)\bmod 1\bigr) < 1$; in particular
\[
\mathrm{ALG} \;\le\; \underbrace{\sum_{(e,t)\in U}\sum_i f_i(e,t)}_{=\;\sum_i \mathrm{dem}(d_i)\cdot k_i^{\mathrm{algo}}} \;+\; |U|.
\]

\paragraph{Bound on $|U|$.}
The edges traversed by HRPES are contained in
\[
E(H_\mathrm{hub}) \;\subseteq\; \underbrace{\bigl(\bigcup_v \pi_v\bigr)}_{\text{access trees}} \;\cup\; \underbrace{\bigl(\bigcup_{s,s'\in S}\pi_{s,s'}\bigr)}_{\text{inter-hub backbone}}.
\]
The access trees together have at most $n-|S|$ edges (Shortest-Path Forest rooted at $S$). The inter-hub backbone has at most $\binom{|S|}{2}\cdot\mathrm{Diam}(G)=O(|S|^2\,\mathrm{Diam}(G))=O(n\log n\cdot\mathrm{Diam}(G))$ edges. Each edge $e\in E(H_\mathrm{hub})$ is used by at most $h$ demands, but each demand contributes at most one time per edge, so
\[
|U|\;\le\;\sum_{e\in E(H_\mathrm{hub})}\bigl|\{\tau_i(e):e\in P_i\}\bigr|\;\le\;\min\bigl\{h\cdot|E(H_\mathrm{hub})|,\;|E(H_\mathrm{hub})|\cdot a\bigr\}.
\]
A tighter bound comes from the optional consolidation: aligning all demands sharing an edge to a common time gives $|U|\le|E(H_\mathrm{hub})|$. Without consolidation:
\[
|U|\;\le\;|E(H_\mathrm{hub})|\cdot \min\{h,a\}\;=\;O\bigl((n+n\log n\cdot\mathrm{Diam}(G))\cdot\min\{h,a\}\bigr).
\]

\begin{theorem}[Bicriteria bound for HRPES]
\label{thm:main}
HRPES outputs a feasible MAL-C solution of cost
\[
\mathrm{ALG} \;\le\; 2\,\mathrm{Diam}(G)\cdot \mathrm{OPT}_\mathrm{fluid} \;+\; |U|
\;\le\; 2\,\mathrm{Diam}(G)\cdot \mathrm{OPT}_\mathrm{MAL\text{-}C} \;+\; |U|,
\]
where $|U|=O\bigl((n+n\log n\cdot\mathrm{Diam}(G))\cdot\min\{h,a\}\bigr)$ in the worst case, and $|U|\le|E(H_\mathrm{hub})|=O(n\log n\cdot\mathrm{Diam}(G))$ when the optional consolidation step is applied.
\end{theorem}
\begin{proof}
Plug the per-demand inequality $k_i^{\mathrm{algo}}\le 2\,\mathrm{Diam}(G)\cdot k_i^*$ into $\sum_i\mathrm{dem}(d_i)k_i^{\mathrm{algo}}\le 2\,\mathrm{Diam}(G)\sum_i\mathrm{dem}(d_i)k_i^*=2\,\mathrm{Diam}(G)\cdot\mathrm{OPT}_\mathrm{fluid}$, then add the ceiling overhead $|U|$. Use $\mathrm{OPT}_\mathrm{fluid}\le\mathrm{OPT}_\mathrm{MAL\text{-}C}$ for the second inequality.
\end{proof}

\subsubsection*{4. What this bound does and does not say}

\begin{enumerate}[noitemsep,leftmargin=*]
  \item \textbf{This is a bicriteria bound, not a multiplicative approximation.}
  Dividing both sides by $\mathrm{OPT}_\mathrm{MAL\text{-}C}$,
  \[
    \frac{\mathrm{ALG}}{\mathrm{OPT}_\mathrm{MAL\text{-}C}} \;\le\; 2\,\mathrm{Diam}(G) \;+\; \frac{|U|}{\mathrm{OPT}_\mathrm{MAL\text{-}C}}.
  \]
  The first term is bounded; the second is bounded only when $\mathrm{OPT}_\mathrm{MAL\text{-}C}=\Omega(|U|)$, i.e.\ in the heavy-traffic regime. In light traffic ($\mathrm{OPT}_\mathrm{MAL\text{-}C}=O(1)$) the ratio can be as large as $\Theta(|U|)$, so the algorithm is \emph{not} an $O(\mathrm{Diam}(G))$-approximation in the worst case.
  \item \textbf{Heavy-traffic asymptotics.}
  Concretely, when $\mathrm{OPT}_\mathrm{MAL\text{-}C}\ge|U|$ the ratio is at most $2\,\mathrm{Diam}(G)+1$.
  \item \textbf{Comparison with CDO25.}
  Carnevale et al.\ (Theorem~3 of \cite{CDO25}, summarised in Table~1 of CM03-1) achieve $O(\sqrt[3]{n\log^2 n\cdot\mathrm{Diam}(G)})$ for $a\ge\lceil 5/3\,\mathrm{Diam}(G)\rceil$ and $O(\sqrt{n\log n})$ for $a\ge\mathrm{Diam}(G)+\lfloor 0.5\,\mathrm{Diam}(G)\rfloor$. HRPES's heavy-traffic ratio $O(\mathrm{Diam}(G))$ improves on these only when $\mathrm{Diam}(G)$ is small (e.g.\ small-world graphs); on graphs with $\mathrm{Diam}(G)=\omega(\sqrt{n\log n})$ it is worse. CDO25 furthermore solves a strictly different objective (number of distinct labels, not vehicle dispatches), so a direct comparison should be made via Lemma~\ref{lem:fluid} rather than by quoting their ratios.
  \item \textbf{Capacity / volume.}
  The MAL-C objective $\sum\lceil\sum_i f_i(e,t)\rceil$ already absorbs unit-vehicle capacity. HRPES does not exploit the freedom to split a demand across multiple paths to fill ceilings; doing so could reduce $|U|$ further but is not analysed here.
\end{enumerate}

\subsubsection*{5. Why this is correct without a global TEG}

\begin{proposition}
HRPES never constructs an object whose size depends on $a$.
\end{proposition}
\begin{proof}
Phase~1 builds $S$, the access forests, and the inter-hub backbone — all in $G$, with size $O(n+n\log n\cdot\mathrm{Diam}(G))$. Phase~2 computes $\tau_i(e)$ for each demand and each edge in $P_i$ in $O(|P_i|)=O(\mathrm{Diam}(G))$ work per demand. Phase~3 emits labels and flows by simple lookup. The DP in Lemma~\ref{lem:fluid} runs over states $(v,k)$ with $k\le n$ — also independent of $a$. No phase enumerates time slots in $\{1,\dots,a\}$.
\end{proof}

\paragraph{Remark on the previous draft.}
An earlier version of this document defined a per-node ``cutoff time'' $t^*(u)=\min_{(u,v)\in H_\mathrm{hub}}\{t^*(v)-w(u,v)\}$ and claimed that demands released after $t^*(\mathrm{src})$ could be handled by a ``mini-TEG of size independent of $a$''. Both claims were incorrect: $t^*(u)$ is not well-defined when $u$ has multiple successors with different deadlines, and the late-release window $a-\mathrm{rel}(d_i)$ can be $\Theta(a)$. The construction above replaces both by per-edge times $\tau_i(e)$, which are well-defined and require no fallback TEG.

\begin{thebibliography}{9}
\bibitem{CDO25} D.~Carnevale, G.~D'Angelo, M.~Olsen. Approximating optimal labelings for temporal connectivity. AAAI 2025.
\end{thebibliography}

\end{document}
